%% kapa: exam=1 round=30 year=2019 period=1 form=A subject=real_estate subject_ko=부동산학원론 seq=35 number=115
%% source: Q-Net 감정평가사 기출문제 (cst00309) artlSeq=5205205 file=extracted/1008/5205205/2019년 제30회 감정평가사 1차시험 문제지/2019년 제30회 감정평가사 1차시험 문제지/2019년도 제30회 감정평가사 1차 1교시 A형.hwp
%% conversion: kordoc 4.24.51 -> tools/gen1.py (auto); review=proofread-sonnet; flags=none
\begin{question}{115}
부동산개발의 개념에 관한 설명으로 옳지 않은 것은?
\begin{choices}
  \item ｢부동산개발업의 관리 및 육성에 관한 법률｣상 부동산개발은 시공을 담당하는 행위를 포함한다.
  \item 부동산개발은 온전하게 운용할 수 있는 부동산을 생산하기 위한 토지와 개량물의 결합이다.
  \item 부동산개발이란 인간에게 생활, 일, 쇼핑, 레저 등의 공간을 제공하기 위한 토지, 노동, 자본 및 기업가적 능력의 결합과정이다.
  \item 부동산개발은 토지조성활동과 건축활동을 포함한다.
  \item 부동산개발은 토지 위에 건물을 지어 이익을 얻기 위해 일정 면적의 토지를 이용하는 과정이다.
\end{choices}
\end{question}
\answer{1}
